\section*{Introduction}

The motivational problem of this course is the following.

\begin{problem}[Shortest Path Problem]
    How do we find the shortest path between two points:
    \begin{enumerate}
        \item in Euclidean Space;
        \item on the Sphere;
        \item on a general curved surface, e.g., potato?
    \end{enumerate}
\end{problem}

In problem (i), we may write the path in the form \(y = y\pqty{x}\), with \(a \leq x \leq b\). Then its length is given by
\[
    L \bqty{y} = \int_{a}^{n} \sqrt{1 + y'^2} \dd{x}.
\]

We can find the length of a parametrised curve similarly in (ii) and (iii).

This is a \emph{functional}, which is a map from a function to a number, and the goal is to find the choice of the function \(y\pqty{x}\) that minimises this functional. We will develop tools in this course to determine minima/maxima  of functionals.